\section{Completion and tangent cones}\label{sec:distances}

We distinguish two path distances until their equality has
been proved. On $M_{\reg}$, let $d_{\reg}$ be the Riemannian
distance of $g^{L^2}$. On all of $M$, call a path admissible
if it is continuous, has a finite partition into smooth paths
lying in single strata on the open subintervals, and has
finite one-sided stratum lengths at the endpoints. Let
$d_{\str}$ be the infimum of these lengths. The same
convention on each finite cone gives its full length distance
$d_{\Cone_E}$. Piecewise smooth paths and their finite-length
limits give the same infimum.

\subsection{Short paths in the regular locus}
The following elementary consequence of semialgebraicity
provides control near the endpoints of transverse perturbations.

\begin{lemma}\label{lem:short-paths}
Let $\Cone$ be one of the cones in
\eqref{eq:offdiag-cone}--\eqref{eq:diagonal-cone}, or a full
normal cone to one of its strata. There is a constant $A_\Cone$
such that any two regular points of radius at most $\rho$
can be joined in the regular locus, within that radial ball,
by a piecewise smooth path of length at most $A_\Cone\rho$.
\end{lemma}

\begin{proof}
Write the full normal representation as $W$ with compact
group $K'$, and let $B^*$ be the module-simple part of the
simultaneous zero set of its real and complex moment maps.
The simple locus of $\Cone$ is a connected complex open
subset of an irreducible cone. For the two original cones,
irreducibility and normality follow from their moduli germs;
homogeneity extends these properties to the whole cone.
The same argument applies to full normal slices, using the
slice theorem. Simplicity agrees with regularity here by
the stable-locus description of Section~\ref{sec:moduli}
and its slice restriction. Existence and density of simple
points in every normal cone were recalled in
Section~\ref{subsec:slice-input}.

The map $B^*\to\Cone_{\reg}$ is proper with connected
$K'$-orbit fibres. The group $K'$ is a product of unitary
groups, with an optional ineffective scalar quotient, and
is connected. Thus $B^*$ is connected. If the zero
representation itself is simple, the effective group is
trivial and the quotient is a flat vector space; the
assertion is then immediate. Otherwise radial scaling
identifies $B^*$ with its unit section times $(0,\infty)$.
It follows that
\[
 B^*\cap\{0<\|w\|<1\}
\]
is a bounded connected semialgebraic set.

We recall why such a set has finite Euclidean inner diameter,
including when it is not closed. A finite semialgebraic
triangulation uses finitely many open simplices. In each
simplex choose an anchor and join a variable point to that
anchor along the image of the straight segment in the
simplex. Join the finitely many anchors by fixed
semialgebraic arcs in the connected set. This gives finitely
many semialgebraic families of paths. Uniform monotonicity
partitions each coordinate of each path into a bounded
number of monotone intervals. Each coordinate varies over
a bounded range, so its total variation, and hence the
Euclidean length, has a common bound. This is the
triangulation and monotonicity argument of
\cite[Sections~2 and~3, especially Theorem~3.4]{vdd98}.
Semialgebraic arcs are piecewise smooth; endpoint derivatives
may be unbounded, but their lengths are finite.

Projection from the balanced locus to its smooth quotient
does not increase lengths. We obtain a uniform bound for
regular points of cone radius less than one. Dilation gives
the bound $A_\Cone\rho$ for an open radial ball. For points
on its boundary, first contract both endpoints by a factor
$1-\delta$, join the contracted points, and append the two
regular radial segments. Their total added length is at
most $2\delta\rho$. This proves the assertion for the
closed ball as well, after increasing the constant.
\end{proof}

\begin{proposition}\label{prop:completion}
The completion of $(M_{\reg},d_{\reg})$ is canonically
homeomorphic to $M$, and its distance equals $d_{\str}$.
In particular it is a compact geodesic space.
\end{proposition}

\begin{proof}
Shrink each fixed chart in Theorem~\ref{thm:main}(ii) so
that its squared metric bounds lie between $1/2$ and $2$.
Lemma~\ref{lem:short-paths} then shows that the $d_{\reg}$
diameter of regular points in a radial ball of radius
$\rho$ is at most $A\rho$. Every sequence of regular
points converging to the centre therefore determines the
same completion point.

There is also a lower bound for leaving a chart. If an
admissible path starts at cone radius $\rho$ and first
reaches the boundary of a concentric chart of radius $R$,
its length before exit is at least
\begin{equation}\label{eq:exit}
 (R-\rho)/\sqrt2.
\end{equation}
This follows by integrating the radial $1$-Lipschitz
inequality for the cone and the coarse tensor lower bound.
In particular, distinct points of $M$ cannot represent
the same completion point. Conversely, a $d_{\reg}$-Cauchy
sequence has an algebraically convergent subsequence because
$M$ is compact. The diameter estimate identifies its limit
in the completion; the Cauchy property gives the same limit
for the full sequence. Thus no additional completion points
occur. These estimates also prove that the completion
topology is the original topology of $M$. Denote this
completed distance by $D$.

We next compare lengths. Let $\gamma:[a,b]\to M$ be a
$C^1$ path contained in a single smooth stratum. Given
$\eta>0$, cover its compact image by finitely many centred
charts whose metric distortion is between $1-\eta$ and
$1+\eta$. In each such chart the cone splits as
$U\times\mathcal N$, and the stratum through its centre
is $U\times\{0\}$. Subdivide $[a,b]$ so that each
subarc is contained in one chart and write it there as
$(u(t),0)$.

Choose regular points $z_j\in\mathcal N$ tending to
the vertex. They exist by the simple-point property in
Section~\ref{subsec:slice-input}. The paths
\begin{equation}\label{eq:transverse-path}
 t\longmapsto (u(t),z_j)
\end{equation}
lie in the regular locus for large $j$. Their lengths in
the exact product are the Euclidean length of $u$.
The two tensor bounds give, for each subarc,
\[
 \limsup_j L_{g^{L^2}}(\Psi(u,z_j))
 \le \sqrt{\frac{1+\eta}{1-\eta}}\,L_{g^{L^2}}(\gamma).
\]
Their endpoints converge in the original topology, hence
in $D$ by the preceding paragraph. Taking endpoint
distance limits and summing with the triangle inequality
gives the same upper bound for $D(\gamma(a),\gamma(b))$.
No common endpoint is required for the different
perturbed subarcs. Letting $\eta$ tend to zero gives
$D(\gamma(a),\gamma(b))\le L(\gamma)$.

For an admissible path whose open pieces have endpoints
in lower strata, truncate each piece at both ends and
apply this argument. The finite one-sided lengths and
continuity of $D$ allow the truncation parameters to
tend to zero. Therefore $D\le d_{\str}$ on $M$.
On regular points, every regular path is admissible, so
$d_{\str}\le d_{\reg}=D$. Finally, $d_{\str}$ is
continuous in the original topology: in a chart centred
at any point, its radial paths have length at most
$\sqrt2\rho$, and the triangle inequality gives
continuity in both arguments. Approximation by regular
endpoints extends the reverse inequality to all of $M$.
Thus $D=d_{\str}$.

The completion of a length space is a length space.
It is compact here, so the usual compactness argument
for almost minimizing paths supplies minimizing
geodesics. This proves the last assertion.
\end{proof}

\subsection{Quantitative convergence to the cone}
We can now transfer the quadratic tensor comparison to
the completed distance. Fix a centre $p=[E]$ and its
chart $\Psi_E$. Competitors joining points of radius
at most $\rho$ can be chosen with length $O(\rho)$,
using radial paths through the centre. A competitor
of this length cannot leave a fixed chart when $\rho$
is small, by \eqref{eq:exit}. Within the chart, the
coarse lower bound further shows that its radius is
$O(\rho)$ throughout. On all such competitors the
tensor distortion is therefore $1+O_E(\rho^2)$.

Apply this reasoning first to almost minimizing cone
paths and then to almost minimizing admissible paths
in $M$. Proposition~\ref{prop:completion} identifies
the latter infimum with $D$. Taking both infima gives
\begin{equation}\label{eq:distance-comparison}
 (1-C'_E\rho^2)d_{\Cone_E}(x,y)
 \le D(\Psi_E x,\Psi_E y)
 \le(1+C'_E\rho^2)d_{\Cone_E}(x,y)
 \quad(r(x),r(y)\le\rho).
\end{equation}
The slightly larger radius used by the competitors
is absorbed into $C'_E$. Setting $y=0$ proves
\eqref{eq:radius-main}.

For completeness, fix $L<\infty$ and put
$x=\epsilon\hat x$, $y=\epsilon\hat y$ in
\eqref{eq:distance-comparison}. On the cone ball of
radius $L$, the map
$\hat x\mapsto\Psi_E(\epsilon\hat x)$ has distance
distortion $O_{E,L}(\epsilon^2)$ when the target
distance is divided by $\epsilon$. The radial formula
shows that the preimage of the closed target unit
ball lies between cone radii $1-O_E(\epsilon^2)$
and $1+O_E(\epsilon^2)$. Radial adjustment to radius
one costs $O_E(\epsilon^2)$ in the rescaled metric.
This gives a correspondence of the two closed unit
balls with that distortion, hence the
Gromov--Hausdorff estimate in
Corollary~\ref{cor:completion-main}. The construction
works for every sequence $\epsilon\downarrow0$;
the metric tangent cone is consequently unique.
